\section{Begrippenlijst}
\label{app:begrippenlijst}

\renewcommand{\arraystretch}{1.4}
\begin{longtable}{p{4.5cm} p{10cm}}
    \rowcolor{green!20}
    \textbf{Begrip} & \textbf{Omschrijving} \\
    \hline
    \endfirsthead

    \rowcolor{green!20}
    \textbf{Begrip} & \textbf{Omschrijving} \\
    \hline
    \endhead

    Accelerometer &
    Sensor die de versnelling van een object meet langs drie assen.
    Meet naast bewegingsversnelling ook altijd de zwaartekracht. \\

    ASM330LHH &
    Automotive-grade IMU van STMicroelectronics met een 3-assige
    accelerometer en 3-assige gyroscoop, geïntegreerd op de
    ProRTK tracker. \\

    Bias &
    Systematische meetafwijking van een sensor die aanwezig is
    ongeacht de werkelijke beweging. Accumuleert bij integratie
    tot een toenemende positiefout. \\

    BDU (Block Data Update) &
    Instelling van de IMU die ervoor zorgt dat de uitvoerregisters
    pas worden bijgewerkt nadat de microcontroller ze volledig
    heeft uitgelezen, om incorrecte meetwaarden te voorkomen. \\

    CAN-bus &
    Communicatieprotocol voor betrouwbare datacommunicatie tussen
    elektronische componenten. \\

    Covariantiematrix ($P$) &
    Matrix in het EKF die bijhoudt hoe zeker het filter is over
    elke grootheid. Grote diagonaalwaarden duiden op grote
    onzekerheid. Neemt toe bij de predictiestap en af bij
    de updatestap. \\

    Complementary filter &
    Eenvoudig filteralgoritme dat twee sensordatastromen combineert
    op basis van een vaste weging. Wordt gebruikt om de
    hellingshoek van de tracker te bepalen met behulp van de
    accelerometer- en gyroscoopdata. \\

    DPS (Degrees Per Second) &
    Eenheid voor hoeksnelheid, gebruikt voor het aangeven van het
    meetbereik van een gyroscoop. \\

    ECEF (Earth-Centered Earth-Fixed) &
    Driedimensionaal coördinatenstelsel met de oorsprong in het
    middelpunt van de aarde, waarbij X, Y en Z in meters worden
    uitgedrukt. Wordt als tussenstap gebruikt bij de omrekening
    van geografische coördinaten naar het ENU-stelsel. \\

    EKF (Extended Kalman Filter) &
    Variant van het Kalman Filter die niet-lineaire systemen kan
    verwerken door het bewegingsmodel bij elke tijdstap lokaal
    te lineariseren. Vormt de kern van het sensorfusie systeem
    in dit project. \\

    ENU (East-North-Up) &
    Lokaal coördinatenstelsel waarbij X naar het oosten wijst,
    Y naar het noorden en Z omhoog. Wordt gebruikt als intern
    coördinatenstelsel van het EKF. \\

    FPU (Floating-Point Unit) &
    Hardware-eenheid in de processor die kommaberekeningen
    uitvoert. Belangrijk voor de intensieve matrixoperaties
    in het EKF. \\

    GNSS (Global Navigation Satellite System) &
    Overkoepelende term voor satellietnavigatiesystemen zoals GPS
    (VS), Galileo (EU), GLONASS (Rusland) en BeiDou (China).
    Bepaalt de positie via trilateratie op basis van
    aankomsttijden van satellietsignalen. \\

    Gyroscoop &
    Sensor die de hoeksnelheid van een object meet langs drie
    assen. Door integratie over de tijd kan de
    oriëntatieverandering worden bepaald. Gevoelig voor drift
    bij langdurig gebruik. \\

    hAcc (Horizontal Accuracy) &
    Schatting van de horizontale positienauwkeurigheid zoals
    vermeld door de u-blox GNSS-ontvanger. Beschikbaar op 25\,Hz. \\

    IMU (Inertial Measurement Unit) &
    Sensor die de beweging van een object meet via traagheid,
    onafhankelijk van externe referenties. Bestaat typisch uit
    een accelerometer en een gyroscoop (soms ook een magnetometer). \\

    Innovatie &
    Het verschil tussen de door het EKF voorspelde positie en de
    werkelijk gemeten RTK positie. Wordt gebruikt om de
    toestandsschatting bij te sturen in de updatestap. \\

    Jacobiaan ($F$) &
    Matrix van partiële afgeleiden die beschrijft hoe een kleine
    verandering in één grootheid invloed heeft op alle andere
    grootheden na één tijdstap. Wordt gebruikt om het
    niet lineaire bewegingsmodel lineair te maken in het ekf. \\

    Kalman gain ($K$) &
    Wegingsfactor in het EKF die dynamisch bepaalt hoeveel gewicht
    de GNSS meting krijgt ten opzichte van de eigen voorspelling (IMU meting)
    van het filter. Groot bij betrouwbare metingen, klein bij
    onzekere metingen. \\

    Multipath &
    Effect waarbij een GNSS signaal reflecteert tegen
    objecten voordat het de antenne bereikt. Veroorzaakt een
    positiefout doordat de ontvanger een langere signaalweg meet
    dan de werkelijke afstand. \\

    NAV-COV &
    UBX bericht van de u-blox ontvanger dat de volledige
    positiecovariantiematrix bevat, met varianties in oost en
    noordrichting. \\

    NAV-PVT &
    UBX bericht van de u-blox ontvanger dat positie, snelheid
    en nauwkeurigheidsschatting (hAcc) bevat. \\

    NIS (Normalized Innovation Squared) &
    Statistische toets die meet hoe verrassend een binnenkomende
    GNSS-meting is in vergelijking met de huidige staat van het filter. Wordt gebruikt
    om grote foutive sprongen te verwerpen. \\

    nRF52832 &
    Microcontroller van Nordic Semiconductor met een ARM
    Cortex-M4F processor en hardware FPU. \\

    ODR (Output Data Rate) &
    De frequentie waarmee een sensor nieuwe meetwaarden levert,
    uitgedrukt in Hz. \\

    Procesruismatrix ($Q$) &
    Matrix in het EKF die per tijdstap modelleert hoeveel
    onzekerheid er wordt toegevoegd doordat het bewegingsmodel
    nooit perfect is. Bepaalt hoe snel het filter openstaat
    voor bijsturing. \\

    ProRTK tracker &
    Compacte sporttimer van MYLAPS die draadloos communiceert
    via het X2 Link protocol en een RTK GNSS-ontvanger combineert
    met een IMU op één printplaat. \\

    RTK (Real-Time Kinematic) &
    Techniek die de nauwkeurigheid van GNSS verbetert tot
    1--2\,cm door gebruik te maken van een base station dat
    real time correcties stuurt naar de mobiele ontvanger. \\

    RTK Fixed &
    Modus van de RTK ontvanger waarbij de
    `carrier phase ambiguity' volledig is opgelost. Geeft de hoogste
    nauwkeurigheid van 1--2\,cm. \\

    RTK Float &
    Operatiemodus van de RTK ontvanger waarbij de
    `carrier phase ambiguity' nog niet definitief is opgelost.
    Nauwkeurigheid typisch 10--50\,cm maar kan ook meer zijn. \\

    Sensorfusie &
    Het combineren van meerdere sensordatastromen tot één
    betrouwbare schatting, waarbij de voordelen van elke sensor
    worden benut en de nadelen worden gecompenseerd. \\

    SPI (Serial Peripheral Interface) &
    Synchroon serieel communicatieprotocol voor communicatie
    tussen de microcontroller en hardware. De IMU,
    seriële flash en CAN controller delen dezelfde SPI bus
    op de ProRTK. \\

    State vector ($\mathbf{x}$) &
    Vector die alle geschatte grootheden van het systeem bevat:
    positie ($X$, $Y$), snelheid ($v$), heading ($\theta$),
    accelerometerbias ($b_a$) en gyroscoop bias ($b_g$). \\

    UART &
    Serieel communicatieprotocol voor communicatie tussen de
    microcontroller en de GNSS ontvanger. \\

    u-blox X20P &
    RTK GNSS ontvanger geïntegreerd in de ProRTK tracker.
    Levert positie, snelheid en nauwkeurigheidsmarges. \\

    Unicycle-model &
    Vereenvoudigd rijmodel uit de robotica dat ervan uitgaat dat
    een object vooruit beweegt met een bepaalde snelheid en
    draaisnelheid, zonder zijwaartse slip. \\

    WGS84 &
    Dit is een wereldwijd coördinatenstelsel dat geografische locaties bepaalt op basis van een wiskundig model van de aarde (de referentie-ellipsoïde). Het is de wereldwijde standaard die door alle GNSS ontvangers (zoals gps) wordt gebruikt. \\

    X2 Link &
    Draadloos communicatieprotocol van MYLAPS voor de overdracht
    van de nuttige paramaters naar de server. Hiermee kunnen deze parameters uitgelezen worden doormiddel van de ontwikkelde applicaties. \\

\end{longtable}